% Options for packages loaded elsewhere
\PassOptionsToPackage{unicode}{hyperref}
\PassOptionsToPackage{hyphens}{url}
\PassOptionsToPackage{dvipsnames,svgnames,x11names}{xcolor}
%
\documentclass[
  11pt,
]{article}
\usepackage{amsmath,amssymb}
\usepackage{iftex}
\ifPDFTeX
  \usepackage[T1]{fontenc}
  \usepackage[utf8]{inputenc}
  \usepackage{textcomp} % provide euro and other symbols
\else % if luatex or xetex
  \usepackage{unicode-math} % this also loads fontspec
  \defaultfontfeatures{Scale=MatchLowercase}
  \defaultfontfeatures[\rmfamily]{Ligatures=TeX,Scale=1}
\fi
\usepackage{lmodern}
\ifPDFTeX\else
  % xetex/luatex font selection
\fi
% Use upquote if available, for straight quotes in verbatim environments
\IfFileExists{upquote.sty}{\usepackage{upquote}}{}
\IfFileExists{microtype.sty}{% use microtype if available
  \usepackage[]{microtype}
  \UseMicrotypeSet[protrusion]{basicmath} % disable protrusion for tt fonts
}{}
\makeatletter
\@ifundefined{KOMAClassName}{% if non-KOMA class
  \IfFileExists{parskip.sty}{%
    \usepackage{parskip}
  }{% else
    \setlength{\parindent}{0pt}
    \setlength{\parskip}{6pt plus 2pt minus 1pt}}
}{% if KOMA class
  \KOMAoptions{parskip=half}}
\makeatother
\usepackage{xcolor}
\usepackage[margin=1in]{geometry}
\setlength{\emergencystretch}{3em} % prevent overfull lines
\providecommand{\tightlist}{%
  \setlength{\itemsep}{0pt}\setlength{\parskip}{0pt}}
\setcounter{secnumdepth}{-\maxdimen} % remove section numbering
\ifLuaTeX
  \usepackage{selnolig}  % disable illegal ligatures
\fi
\IfFileExists{bookmark.sty}{\usepackage{bookmark}}{\usepackage{hyperref}}
\IfFileExists{xurl.sty}{\usepackage{xurl}}{} % add URL line breaks if available
\urlstyle{same}
\hypersetup{
  pdftitle={Three framework-specification questions to unblock first-principles (r\_e) derivation},
  pdfauthor={Trey Morris with Claude Opus 4.7},
  pdfsubject={Three framework-specification questions for hypothesis-(i)
first-principles r\_e derivation (follow-up to candidates note +
overnight progress)},
  colorlinks=true,
  linkcolor={blue},
  filecolor={Maroon},
  citecolor={Blue},
  urlcolor={blue},
  pdfcreator={LaTeX via pandoc}}

\title{Three framework-specification questions to unblock
first-principles (r\_e) derivation}
\author{Trey Morris with Claude Opus 4.7}
\date{2026-05-27}

\begin{document}
\maketitle

\hypertarget{three-framework-specification-questions-for-first-principles-r_e-derivation-for-tepper-gill}{%
\section{\texorpdfstring{Three framework-specification questions for
first-principles \(r_e\) derivation --- for Tepper
Gill}{Three framework-specification questions for first-principles r\_e derivation --- for Tepper Gill}}\label{three-framework-specification-questions-for-first-principles-r_e-derivation-for-tepper-gill}}

\textbf{Date:} 2026-05-26 (PM, after the day's overnight progress runs).
\textbf{From:} Trey Morris (with Claude Opus 4.7). \textbf{Re:}
Follow-up to the
\href{2026-05_re_derivation_candidates_for_gill.md}{candidates note} and
the \href{2026-05_re_triangulation_followup_for_gill.md}{triangulation
note}. Tracks against
\href{https://github.com/temoTxt/PyPhysics/issues/75}{issue \#75} under
master \href{https://github.com/temoTxt/PyPhysics/issues/67}{\#67}.

\begin{center}\rule{0.5\linewidth}{0.5pt}\end{center}

Tepper --- following the candidates note we sent earlier today (the PM
revision adds overnight progress from three parallel research branches),
I want to consolidate the three framework-specification questions that
surfaced as load-bearing across all three candidate branches. Each of
these would unblock a first-principles derivation of \(r_e/r_0\) that
goes beyond the algebraic inheritance from QED's \(a_e(\alpha)\) that
the current closed-form result delivers.

\hypertarget{what-the-three-branches-converged-on}{%
\subsection{What the three branches converged
on}\label{what-the-three-branches-converged-on}}

\begin{itemize}
\item
  \textbf{Candidate 3
  (\href{https://github.com/temoTxt/PyPhysics/issues/66}{\#66},
  \href{https://github.com/temoTxt/PyPhysics/pull/70}{PR \#70})}
  confirmed via a line-by-line read of §III.D Eqs. (III.18)--(III.23)
  that the as-published apparatus leaves \(r_e/r_0\) as a free parameter
  --- the algebraic-back-fit match to the closed-form
  \(r_e/r_0 = (2 - \alpha/(2\pi))/(4 + \alpha/\pi)\) is necessary but
  not sufficient for intentional Schwinger encoding. Iter-5 of that
  branch added a cross-particle test against the FNAL Muon \(g-2\) 2023
  measurement: \(r_\mu/r_0^\mu - r_e/r_0^e = 3.13\times 10^{-6}\) at
  \(\sim 57{,}000\,\sigma_{a_\mu}\), which \textbf{empirically excludes
  any universal closed-form \(r/r_0 = f(\alpha)\) in \(\alpha\) alone}.
\item
  \textbf{Candidate 2
  (\href{https://github.com/temoTxt/PyPhysics/issues/65}{\#65},
  \href{https://github.com/temoTxt/PyPhysics/pull/69}{PR \#69})} found
  that the published expanded \(K_D\) of (III.4) is structurally
  inadequate to pin \(r_e\) via the variational principle: no scale
  exists where both (a) the NR expansion is valid AND (b) the cutoff
  couples meaningfully to the closure equation. At the electron-radius
  scale (\(\hat{r}_e \sim 1\)), kinetic energy dominates by four orders,
  so the NR expansion is invalid; at the Bohr scale where the expansion
  is valid, the cutoff \(\hat{r}_e\) is suppressed by \(\alpha^2\) and
  the closure matches plain hydrogen to 10 sig figs.
\item
  \textbf{Candidate 1
  (\href{https://github.com/temoTxt/PyPhysics/issues/64}{\#64},
  \href{https://github.com/temoTxt/PyPhysics/pull/72}{PR \#72})}
  delivered the closed form \(r_e/r_0 = (2 - a_e)/(2(2 + a_e))\)
  matching the triangulated value to \(3.45\times 10^{-13}\). The closed
  form is the algebraic inverse of the framework's existing g-formula
  (Eq. III.22) evaluated at QED's \(a_e(\alpha)\) --- i.e.,
  \emph{reproduction-by-inheritance} under what we have been calling
  hypothesis (ii) (photon propagator unchanged, all dual structure
  absorbed into the (II.3) Pauli source). The framework does not
  independently re-derive Schwinger's vertex result; it inherits it.
\end{itemize}

The convergent gap across all three is \textbf{a first-principles
derivation of \(a_e\) from framework-internal dynamics}, which requires
three framework specifications the published apparatus does not
currently supply. We are not asking you to do the derivation --- we are
asking whether the framework specifies, or has implicit choices for,
these three pieces. Each is small enough to be answered in a few
sentences.

\hypertarget{the-three-questions}{%
\subsection{The three questions}\label{the-three-questions}}

\hypertarget{the-proper-time-photon-propagator-d_ftaux-y}{%
\subsubsection{\texorpdfstring{1. The proper-time photon propagator
\(D_F^{(\tau)}(x-y)\)}{1. The proper-time photon propagator D\_F\^{}\{(\textbackslash tau)\}(x-y)}}\label{the-proper-time-photon-propagator-d_ftaux-y}}

Two natural readings emerge from the framework's existing structure:

\begin{itemize}
\tightlist
\item
  \textbf{(i-a) Schwinger proper-time with \(b\)-dispersion}:
  \(k^2 = (\omega/b)^2 - \mathbf{k}^2\), with
  \(b = \sqrt{c^2 + \mathbf{u}^2}\). The photon dispersion inherits
  \(b\)-factors from the dual-current structure
  (\(J^\mu_{\rm Gill} = (b\rho, \mathbf{J})\), Maxwell paper Eq. 12).
\item
  \textbf{(i-b) Standard Feynman propagator with \(b/c\) conversion
  absorbed into the source's coupling} rather than the propagator:
  \(k^2 = (\omega/c)^2 - \mathbf{k}^2\), with the dual structure
  entering through modified vertex factors.
\end{itemize}

The two give different numerical \(r_e/r_0\) at the same order in
\(\alpha\) --- at the \(\log(b/c)\) level, discriminable at the
framework's precision. Which (if either) is the framework's intended
choice?

\hypertarget{the-bound-state-propagator-g_cxye-in-the-ii.3-potential-in-the-mass-kernel}{%
\subsubsection{\texorpdfstring{2. The bound-state propagator
\(G_C(x,y;E)\) in the (II.3) ``potential-in-the-mass''
kernel}{2. The bound-state propagator G\_C(x,y;E) in the (II.3) ``potential-in-the-mass'' kernel}}\label{the-bound-state-propagator-g_cxye-in-the-ii.3-potential-in-the-mass-kernel}}

For the bound-electron one-loop self-energy, we need a propagator that
handles the Coulomb potential non-perturbatively at the cutoff scale,
where the NR expansion of (III.4) is invalid (per Candidate 2's iter-4
diagnostic above). Standard QED uses the Furry-picture Coulomb-bound
propagator, summing the external Coulomb potential to all orders in
\(Z\alpha\) before perturbing in \(\alpha\). Does the framework specify
(or implicitly use) a dual analog of the Furry-picture construction, or
some other treatment of the bound-state propagator under the (II.3)
kernel?

\hypertarget{the-mass-renormalisation-prescription-at-the-cutoff}{%
\subsubsection{3. The mass-renormalisation prescription at the
cutoff}\label{the-mass-renormalisation-prescription-at-the-cutoff}}

The natural framework-internal closure is

\[\langle K_D\rangle_{r_e} \;=\; m_e c^2 \;+\; \Delta E_{\rm bind}^{\rm framework} \;+\; \Delta E_{\rm SE}^{\rm framework}.\]

The textbook NR working hypothesis
(\(\Delta E_{\rm bind} = \langle V_0\rangle\),
\(\Delta E_{\rm SE} = 0\)) gives a trivially-zero closure under the
published expanded \(K_D\) (per Candidate 2's iter-4 result above). What
is the framework's intended form for
\(\Delta E_{\rm SE}^{\rm framework}(r_e)\) at the cutoff scale --- the
analog of QED's on-shell mass renormalisation condition
\(m_{\rm phys} = m + \Sigma(p)|_{p^2 = m^2c^2}\) in the dual framework's
proper-time formulation?

\hypertarget{posture}{%
\subsection{Posture}\label{posture}}

Any one of these would let us advance the corresponding candidate; all
three would let us attempt the full hypothesis-(i) calculation. The
thread is \textbf{non-urgent} in the sense that the triangulated
\(r_e/r_0 = 0.499\,420\,509\,912\,831\,7\) already serves as the
campaign's current-best-refinement at the framework's precision floor
and is consistent with the framework's apparatus at every observable in
the joint fit. The substantive content of the questions is whether the
framework's apparatus, as you intend it, \emph{contains} the
specifications that would close the gap from ``empirical cutoff'' to
``framework-derived quantity'' --- or whether closing that gap is
genuinely new theory work that the published apparatus does not yet
support.

If the answer to any of (1), (2), (3) is ``the framework does not yet
specify this,'' that is itself a useful finding: it tells us the
campaign's terminal verdict on \(r_e/r_0\) is \textbf{Branch B with
QED-inheritance, pending future framework development} rather than
\textbf{Branch A from first principles}. Either outcome closes the
question honestly; the questions here are about which honest framing is
the right one.

We will record your responses on
\href{https://github.com/temoTxt/PyPhysics/issues/75}{issue \#75} and
pick up the derivation work in a follow-on branch once the
specifications are in hand.

--- Trey

\end{document}
